% !TEX root = ../../main-jos.tex
\section{结束语}

本文面向微服务可观测性场景，提出分布式定向日志采集框架与原型组件，主要贡献体现在以下三个方面：

（1）\textbf{网关驱动的动态 Top-$K$ 关注清单生成}。以网关流量日志作为策略源，通过高价值过滤、URL 聚类泛化与 Top-$K$ 筛选（算法~\ref{alg:attention}），动态识别异常与慢请求 URL 模式，生成关注清单经 Redis 下发至各节点。该策略源创新使得定向采集决策由流量驱动而非静态配置。

（2）\textbf{跨三层的定向策略转换算法}。通过三次转换贯通网关预筛选、节点定向采集与中心二次过滤，实现策略在不同层级的语义一致（图~\ref{fig:transform}），避免了单层采集器配置的局限性。

（3）\textbf{资源受限下的固定缓存与压力感知退避机制}。在 Sidecar 中引入固定上限缓存块与压力感知指数退避，辅以 BoltDB 本地兜底，保障内存有界传输，适应边缘资源受限环境。

基于 Go、gRPC、Redis 与 Loki 的原型在 WSL/Docker 八节点集群上完成端到端验证，并通过 16/32 节点规模复核与基于真实分布的多进程模拟器补充了同频率对照、Promtail 静态过滤复现与组件消融。相较同架构全量采集，定向模式将 Loki 入库量控制在百条量级（八节点实测 72 条 vs.\ 全量 4388$\pm$1 条，详见表~\ref{tab:compare}）；该降幅为策略过滤与源发射频率差异的联合效果，其中策略过滤在同频率模拟下独立降幅为 67.8\%（§6.3）。清空统计后的规模复核显示，16 节点网关流量日志增量为 6703.33$\pm$91.78 条/轮、定向模式 Loki 入库为 48.0$\pm$5.2 条/轮，32 节点网关流量日志增量为 6692.67$\pm$121.88 条/轮、定向模式 Loki 入库为 99.0$\pm$0.0 条/轮；32 节点下关注清单达到 Top-$K{=}50$ 上限。Promtail 静态过滤复现中，本文定向策略入库 233.3 条，低于静态规则的 276.3 条；组件消融中，关闭关注清单使入库量增加 211.7\%，进一步说明动态清单是采集前端减量的关键环节。Agent CPU 绝对值均低于 0.1\%。从 AIOps 生态看，本文工作补充了标准化框架中的数据采集前端能力；从日志故障诊断链路看，本文降低了后端解析、异常检测与根因定位的输入规模。

未来工作按优先级组织：近期将在 Kubernetes 环境开展 64 节点云原生实测，并补充 Top-$K$ 参数敏感性分析；中期将引入 DeathStarBench、Alibaba 生产迹等真实微服务负载验证长尾分布覆盖率，并构建成本量化模型将入库降幅转化为实际成本节约；远期将探索与 OpenTelemetry 追踪上下文、eBPF 应用日志挂钩的集成，以及生产级 gRPC TLS、多租户策略隔离与关注清单 JSON 规范的开源与标准化对接。
